\begin{abstract}
El diagnóstico temprano de enfermedades en los cultivos depende de la inspección visual de
especialistas, un recurso escaso para muchos productores. En este trabajo desarrollamos un
prototipo de clasificación de 15 clases de hojas sanas y enfermas de pimiento, papa y tomate
del dataset PlantVillage, basado en arquitecturas livianas preentrenadas en ImageNet y
entrenado íntegramente en CPU. Para evitar la fuga de datos, particionamos el dataset de forma
estratificada y agrupada por imágenes casi duplicadas. Comparamos cuatro backbones mediante
linear probe y aplicamos fine-tuning parcial a MobileNetV3-Small y EfficientNet-B0, con y sin
data augmentation. EfficientNet-B0 con augmentation alcanza un macro-F1 de 0,993 en test, con
una latencia de 17,5\,ms por imagen en CPU. El augmentation no modifica el desempeño en el test
original, pero es decisivo ante perturbaciones: con desenfoque, el macro-F1 cae 0,23 sin
augmentation y solo 0,013 con él. Un análisis cuantitativo con Grad-CAM muestra que el 65,7\,\%
de la activación recae sobre la hoja, que ocupa el 40,2\,\% de la imagen, y que los errores se
deben a síntomas ambiguos más que a una atención desviada hacia el fondo. Sin embargo, sobre
imágenes de campo del dataset PlantDoc el accuracy cae al 36,7\,\%. Estos resultados indican
que un desempeño cercano al 100\,\% en PlantVillage no garantiza la utilidad del modelo en
condiciones reales.
\end{abstract}

\begin{IEEEkeywords}
clasificación de imágenes, enfermedades de plantas, transfer learning, PlantVillage, EfficientNet, MobileNetV3, Grad-CAM
\end{IEEEkeywords}
